% ------------------------------------------------
% FUENTES
% ------------------------------------------------
\usepackage{fontspec}

\setmainfont{Times New Roman}[
  Ligatures=TeX,
  AutoFakeSmallCaps=true
]

\setmonofont{Courier New}

% ------------------------------------------------
% MÁRGENES Y GEOMETRÍA
% ------------------------------------------------
\usepackage[
  left=2.5cm,
  right=2.5cm,
  top=2.5cm,
  bottom=2.5cm
]{geometry}

\setlength{\headheight}{15pt}

% ------------------------------------------------
% INTERLINEADO Y PÁRRAFOS
% ------------------------------------------------
\usepackage{setspace}
\onehalfspacing
\raggedbottom

\setlength{\parindent}{1.25cm}
\setlength{\parskip}{10pt}

% ------------------------------------------------
% TABLAS E IMÁGENES
% ------------------------------------------------
\usepackage{graphicx}
\usepackage{caption}
\usepackage{float}
\usepackage{longtable}
\usepackage{array}
\usepackage{booktabs}
\usepackage[table]{xcolor}
\usepackage{tabularx}
\usepackage{multirow}
\usepackage{ragged2e}
\usepackage[most]{tcolorbox}
\usepackage{pgfplots}
\pgfplotsset{compat=1.18}
\usepgfplotslibrary{groupplots}

\definecolor{lightgray}{gray}{0.8}
\definecolor{verylightgray}{gray}{0.9}

% Azul corporativo
\definecolor{umaBlue}{cmyk}{1.00, 0.70, 0.00, 0.10}
\definecolor{umaCeleste}{cmyk}{0.80, 0.17, 0.00, 0.14}

\newcolumntype{L}{>{\RaggedRight\arraybackslash}X}
\newcolumntype{C}{>{\Centering\arraybackslash}X}

% --- DEFINICIÓN DE DATOS DEL TRABAJO ---
\title{WebQuest: Criptografía RSA} % <--- Escribe aquí tu título real
\author{Alejandro Muñoz Azaustre}

% Esto crea los comandos \elTitulo y \elAutor para poder usarlos después
\makeatletter
\let\elTitulo\@title
\let\elAutor\@author
\makeatother
% ------------------------------------------------
% TÍTULOS
% ------------------------------------------------
\usepackage{titlesec}

% ======================================
% ENCABEZADOS Y PIES
% ======================================
\usepackage{fancyhdr}
\usepackage{emptypage}
\usepackage{tikz}
\usetikzlibrary{arrows.meta}
\fancypagestyle{tfg}{%
    \fancyhf{}
    \renewcommand{\headrulewidth}{0pt}
    \renewcommand{\footrulewidth}{0pt}

    \fancyhead[LE]{{\normalsize\color{umaBlue}\bfseries\thepage}\quad
                  \footnotesize\MakeUppercase{\leftmark}}

    \fancyhead[RO]{\footnotesize\MakeUppercase{\rightmark}%
                    \quad{\normalsize\bfseries\color{umaBlue}\thepage}}%
}

\renewcommand{\chaptermark}[1]{\markboth{\thechapter.\space #1}{\thechapter.\space #1}}
\renewcommand{\sectionmark}[1]{\markright{\thesection.\space #1}}

\pagestyle{tfg}

\fancypagestyle{plain}{%
    \fancyhf{}
    \renewcommand{\headrulewidth}{0pt}
    \renewcommand{\footrulewidth}{0pt}
    \fancyfoot[C]{\normalsize\bfseries\color{umaBlue}\thepage}
}

% ======================================
% DISEÑO DE CAPÍTULOS
% ======================================
\titleformat{\chapter}[display]
  {\bfseries\large}
  {\hfill\color{umaBlue}\MakeUppercase{\chaptertitlename}\ \large\thechapter}
  {0.5ex}
  {\color{umaBlue}\titlerule\vspace{1.5ex}\filleft\huge\color{umaBlue}}
  [\vspace{0.5ex}\color{umaBlue}\titlerule]

\titlespacing*{\chapter}{0pt}{-30pt}{30pt}

% ======================================
% DISEÑO DE SECCIONES Y SUBSECCIONES
% ======================================
\titleformat{\section}
  {\normalfont\Large\bfseries\color{umaBlue}}
  {\thesection}
  {1em}
  {}

\titleformat{\subsection}
  {\normalfont\large\bfseries\color{umaBlue}}
  {\thesubsection}
  {1em}
  {}

% ------------------------------------------------
% BIBLIOGRAFÍA
% ------------------------------------------------
\usepackage{csquotes}
\usepackage[style=apa,backend=biber]{biblatex}

\makeatletter
\DefineBibliographyExtras{spanish}{
  \setcounter{smartand}{1}
  \let\lbx@finalnamedelim=\lbx@es@smartand
  \let\lbx@finallistdelim=\lbx@es@smartand
}
\makeatother

% ------------------------------------------------
% COMANDOS PERSONALIZADOS
% ------------------------------------------------
\usepackage{comment}
\usepackage{orcidlink}
\usepackage{ccicons}

\newcommand{\ub}[1]{{\color{umaBlue}\bfseries #1}}
\newcommand{\w}[1]{\textcolor{white}{#1}}
\newcommand{\wb}[1]{\textbf{\textcolor{white}{#1}}}

\newcommand{\capitulo}[2][]{%
  \chapter*{#2}%
  \phantomsection
  \markboth{#2}{#2}%
  \addcontentsline{toc}{chapter}{#2}%
  \if\relax\detokenize{#1}\relax
  \else
    \label{#1}%
  \fi
}

\newcommand{\paginaderecha}{%
  \clearpage
  \ifodd\value{page}%
  \else
    \hbox{}
    \newpage
  \fi
}

% ------------------------------------------------
% CAJAS TIPO LIBRO
% ------------------------------------------------
\tcbset{
  enhanced,
  breakable,
  boxrule=0.7pt,
  arc=2mm,
  left=6pt,
  right=6pt,
  top=6pt,
  bottom=6pt
}

\newtcbtheorem[number within=section]{definicion}{Definición}{
  colback=blue!3,
  colframe=blue!60!black,
  coltitle=black,
  fonttitle=\bfseries,
  boxed title style={
    colback=blue!15,
    colframe=blue!60!black
  },
  attach boxed title to top left={yshift=-2mm,xshift=2mm}
}{def}

\newtcbtheorem[number within=section]{ejemplo}{Ejemplo}{
  colback=green!3,
  colframe=green!45!black,
  coltitle=black,
  fonttitle=\bfseries,
  boxed title style={
    colback=green!12,
    colframe=green!45!black
  },
  attach boxed title to top left={yshift=-2mm,xshift=2mm}
}{ej}

\newtcbtheorem[number within=section]{ejercicio}{Ejercicio}{
  colback=orange!4,
  colframe=orange!70!black,
  coltitle=black,
  fonttitle=\bfseries,
  boxed title style={
    colback=orange!15,
    colframe=orange!70!black
  },
  attach boxed title to top left={yshift=-2mm,xshift=2mm}
}{ejer}

\newtcbtheorem[number within=section]{observacion}{Observación}{
  colback=gray!8,
  colframe=gray!65!black,
  coltitle=black,
  fonttitle=\bfseries,
  boxed title style={
    colback=gray!20,
    colframe=gray!65!black
  },
  attach boxed title to top left={yshift=-2mm,xshift=2mm}
}{obs}

\newtcolorbox{solucion}{
  breakable,
  colback=yellow!5,
  colframe=yellow!50!black,
  title=\textbf{Solución},
  fonttitle=\bfseries
}

% ------------------------------------------------
% HIPERVÍNCULOS: Quarto ya carga hyperref
% ------------------------------------------------
\AtBeginDocument{%
  \hypersetup{
    citecolor=umaBlue,
    filecolor=umaBlue,
    linkcolor=umaBlue,
    urlcolor=umaCeleste
  }
}
 % ------------------------------------------------
% DIAGRAMAS DE RELACIONES / FUNCIONES
% ------------------------------------------------
\usepackage{tikz}
\usetikzlibrary{arrows.meta}

\tikzset{
  relabox/.style={
    draw=umaBlue,
    rounded corners=10pt,
    line width=0.8pt
  },
  reladot/.style={
    circle,
    fill=black,
    inner sep=1.6pt
  },
  relaarrow/.style={
    -{Stealth[length=2.3mm,width=1.8mm]},
    line width=0.9pt,
    draw=umaBlue
  },
  relalabel/.style={
    font=\bfseries\color{umaBlue}
  }
}

% Macro general para representar relaciones A -> B
% #1 = nombre de la relación (f, g, m, ...)
% #2 = flechas \draw ...
\newcommand{\RelacionAB}[2]{%
\begin{tikzpicture}[x=1cm,y=1cm]
    % Etiquetas
    \node[relalabel] at (0,1.85) {X};
    \node[relalabel] at (3.8,1.85) {Y};
    \node[relalabel] at (1.9,1.42) {$#1$};

    % Conjuntos
    \draw[relabox] (-0.5,1.5) rectangle (0.5,-0.25);
    \draw[relabox] (3.3,1.5) rectangle (4.3,-0.25);

    % Elementos de A
    \node[reladot] (A1) at (0,1.15) {};
    \node[reladot] (A2) at (0,0.60) {};
    \node[reladot] (A3) at (0,0.05) {};

    % Elementos de B
    \node[reladot] (B1) at (3.8,1.15) {};
    \node[reladot] (B2) at (3.8,0.60) {};
    \node[reladot] (B3) at (3.8,0.05) {};

    % Flechas personalizadas
    #2
\end{tikzpicture}%
} 
